% ======================================================================
%  Wenhan Cao - Curriculum Vitae
%  Compile with XeLaTeX (twice, for hyperlinks).
% ----------------------------------------------------------------------
%  Structure: each section uses small helper macros defined below, so
%  adding a publication = copying one \pubitem{...}{...}{...}{...} block.
% ======================================================================
\documentclass[11pt,a4paper]{article}

% ---------- packages ----------
\usepackage[a4paper,margin=0.63in,top=0.7in,bottom=0.7in]{geometry}
\usepackage[fontset=fandol]{ctex}
\setmainfont{Latin Modern Roman}
\usepackage{enumitem}
\usepackage{titlesec}
\usepackage[hidelinks]{hyperref}
\usepackage{xcolor}
\hypersetup{colorlinks=true, urlcolor=NavyBlue, linkcolor=NavyBlue}
\definecolor{NavyBlue}{RGB}{0,80,160}
\usepackage{needspace}
\usepackage{setspace}
\usepackage{ragged2e}
\usepackage{etoolbox}
\pagestyle{empty}
\setlength{\parindent}{0pt}
\setlength{\emergencystretch}{1em}

% ================= PAGE MANAGEMENT =================
% (a) Fill each page: \raggedbottom lets TeX stretch/shrink the elastic
%     ("plus/minus") spacing defined below so every page bottom aligns.
\flushbottom
% (b) Never leave a section heading alone at the bottom of a page:
%     require at least 4 lines of room before starting a section.
\pretocmd{\section}{\Needspace*{4\baselineskip}}{}{}
% (c) Density knob: one number controls how tight the CV is.
%     0.8 = compact, 1.0 = normal, 1.3 = airy. Recompile after changing.
\newcommand{\density}{1.0}
% (d) Global line spacing: the biggest lever for filling the last page.
%     1.00 = tight, ~1.10 = airy. Tune together with the font size (10/11pt).
\setstretch{1.04}
\newlength{\gapS}\setlength{\gapS}{4pt plus 3pt minus 1pt}  % between list items
\newlength{\gapM}\setlength{\gapM}{6pt plus 4pt minus 2pt}  % between entries
\setlength{\gapS}{\density\gapS}\setlength{\gapM}{\density\gapM}
% ====================================================

% ---------- section heading style ----------
\titleformat{\section}{\large\bfseries}{}{0pt}{}[\vspace{-0.6em}\rule{\textwidth}{0.5pt}]
\titlespacing*{\section}{0pt}{1.1em plus 1.2em minus 0.4em}{0.6em plus 0.3em minus 0.2em}

% ---------- helper macros ----------
% \entry{Institution}{Date}{Role (italic)}{Place (italic)}
\newcommand{\entry}[4]{%
  \begin{samepage}%
  \textbf{#1} \hfill #2\par\nopagebreak
  {\itshape #3} \hfill {\itshape #4}\par\nopagebreak
  \end{samepage}}

% \pubitem{Authors}{Title}{Venue}{Links}   (leave a field empty if unused)
\newcommand{\pubitem}[4]{%
  \item #1. \textit{#2}. #3.\ifx\relax#4\relax\else\ #4\fi}

% \me  -> your bolded name in author lists
\newcommand{\me}{\textbf{Wenhan Cao}}
\newcommand{\mestar}{\textbf{Wenhan Cao}$^*$}

% Translate shared publication link labels only in the Chinese CV.
\csdef{cvlink@Paper}{论文}
\csdef{cvlink@Code}{代码}
\csdef{cvlink@Poster}{海报}
\csdef{cvlink@Slides}{幻灯片}
\csdef{cvlink@Recording}{报告录像}
\csdef{cvlink@Project}{项目主页}
\newcommand{\plink}[2]{[\href{#2}{\ifcsdef{cvlink@#1}{\csuse{cvlink@#1}}{#1}}]}

% award line:  \award{Text}{Year}
\newcommand{\award}[2]{#1 \hfill #2\par\vspace{2pt plus 2pt}}

% publication list environment (hanging indent, tight spacing)
% numbered list for publications: [1], [2], ...
\newenvironment{publist}[1][{[\arabic*]}]{%
  \begin{enumerate}[leftmargin=2.4em,labelsep=0.6em,label={#1},align=left,
                    itemsep=\gapS,parsep=0pt,topsep=3pt plus 1pt]%
  \interlinepenalty=10000\relax}% no page break inside a single reference
  {\end{enumerate}}

% ======================================================================
\newcommand{\acceptedpending}{已接收}
\newcommand{\earlyaccess}{在线优先出版}
\newcommand{\iclrstatus}{投稿至 ICLR 2027，审稿中}
\newcommand{\tmecstatus}{投稿至 IEEE/ASME Transactions on Mechatronics (\textbf{T-MEC}), 2026，审稿中}
\newcommand{\automaticastatus}{投稿至 Automatica，审稿中}
\begin{document}
\begin{center}
  {\Huge\bfseries 曹文涵}\\[5pt]
  {\large Wenhan Cao}\\[6pt]
  \href{https://www.wenhancao.com}{www.wenhancao.com} \;|\;
  \href{mailto:wenhan_cao@outlook.com}{wenhan\_cao@outlook.com} \;|\;
  \href{https://scholar.google.com/citations?user=43xAy7MAAAAJ&hl=en}{谷歌学术}
\end{center}

\section{研究方向}
面向机器人与自动驾驶等具身智能系统在复杂环境下的可靠运行需求，研究不确定性条件下的学习、状态估计与控制方法。以概率建模与推断为基础，研究感知误差、模型偏差及环境变化对系统性能的影响，探索学习、估计与控制的协同设计，提升自主系统的环境感知、决策与闭环控制能力。长期目标是建立\textbf{物理 AI 的概率理论与方法体系}，为自主系统在开放环境中的安全可靠运行提供理论支撑。

\section{工作经历}
\entry{新加坡国立大学}{2025 年 8 月 -- 至今}
      {Eric and Wendy Schmidt AI in Science 博士后研究员（每年遴选不超过五人）\\电子与计算机工程系}{新加坡}
合作导师：\href{https://sites.google.com/view/lzhao}{赵琳教授}

\section{教育与访问经历}
\entry{清华大学}{2019 年 9 月 -- 2025 年 6 月}
      {博士，车辆与运载学院}{中国北京}
导师：\href{https://scholar.google.com/citations?user=Dxiw1K8AAAAJ&hl=en}{李升波教授}（\textbf{车辆与运载学院党委书记}）、\href{http://www2.coe.pku.edu.cn/faculty/liuchang/}{刘畅教授}\par
科研团队：\href{https://www.svm.tsinghua.edu.cn/essay/4/2083.html}{\textbf{李克强院士团队}}
\vspace{\gapM}

\entry{曼彻斯特大学}{2023 年 1 月 -- 2024 年 7 月}
      {计算机科学方向，访问博士研究生}{英国曼彻斯特}
合作导师：\href{https://panweihit.github.io/}{潘为教授}
\vspace{\gapM}

\entry{慕尼黑工业大学}{2023 年 9 月 -- 2023 年 12 月}
      {优化与控制方向，研究访问}{德国慕尼黑}
合作导师：\href{https://www.ce.cit.tum.de/en/itr/hirche/}{Sandra Hirche 教授}
\vspace{\gapM}

\entry{北京交通大学}{2015 年 9 月 -- 2019 年 6 月}
      {工学学士，电气工程}{中国北京}
成绩排名：1/305

\section{代表性论文}
{\small $^*$ 表示共同贡献。论文题目及作者姓名保留英文原文。完整列表见
\href{https://www.wenhancao.com/data/Publication_list_by_type.pdf}{论文列表}或
\href{https://scholar.google.com/citations?hl=en&user=43xAy7MAAAAJ}{谷歌学术}。}
\begin{publist}

\item \me, Xuyang Chen, Shuyuan Wang \& Lin Zhao. \textit{How Much Information is Needed for Accurate Kalman Filtering?} In Advances in Neural Information Processing Systems (\textbf{NeurIPS}), 2026. \plink{Code}{https://anonymous.4open.science/r/rdkf-4B31}

\pubitem{Shiqi Liu$^*$, \mestar, Chang Liu, Zeyu He, Tianyi Zhang, Yinuo Wang \& Shengbo Eben Li}
        {One Filters All: A Generalist Filter For State Estimation}
        {In Advances in Neural Information Processing Systems (\textbf{NeurIPS}), 38 (Main Conference):21276--21302, 2025}
        {\plink{Paper}{https://openreview.net/forum?id=EGK487IYAW}
         \plink{Poster}{https://neurips.cc/virtual/2025/poster/119149}}

\pubitem{\me, Alexandre Capone, Rishabh Dev Yadav, Sandra Hirche \& Wei Pan}
        {Computation-Aware Learning for Stable Control with Gaussian Process}
        {In Robotics: Science and Systems (\textbf{RSS}), 2024}
        {\plink{Paper}{https://www.roboticsproceedings.org/rss20/p004.html}
         \plink{Poster}{https://drive.google.com/file/d/1rSKLy-SEqgMR7Z8Ius9dhHuqOROflWe0/view?usp=sharing}
         \plink{Recording}{https://www.youtube.com/watch?v=AE0iPGlVM30}}

\pubitem{\me\ \& Wei Pan}
        {Impact of Computation in Integral Reinforcement Learning for Continuous-Time Control}
        {In International Conference on Learning Representations (\textbf{ICLR}), 2024. \textbf{(Spotlight)}}
        {\plink{Paper}{https://openreview.net/forum?id=xJEd8PkdNz}
         \plink{Poster}{https://drive.google.com/file/d/16z45LOEU8xaPaEoFFIc2hfIneyXCCio_/view?usp=sharing}
         \plink{Code}{https://github.com/anonymity678/Computation-Impacts-Control}}

\pubitem{\me, Tianyi Zhang, Zeju Sun, Chang Liu, Stephen S.-T. Yau \& Shengbo Eben Li}
        {Nonlinear Bayesian Filtering with Natural Gradient Gaussian Approximation}
        {IEEE Transactions on Pattern Analysis and Machine Intelligence (\textbf{TPAMI}), 48(7):8581--8597, 2026}
        {\plink{Paper}{https://arxiv.org/abs/2410.15832}
         \plink{Code}{https://github.com/wenhancao/NANO-idlab}
         \plink{Slides}{https://drive.google.com/file/d/1O8qNHLxzhbikJ_4Ae-LQPQEzaFJre81U/view?usp=sharing}}

\pubitem{\me, Shiqi Liu, Chang Liu, Zeyu He, Stephen S.-T. Yau \& Shengbo Eben Li}
        {Convolutional Bayesian Filtering}
        {IEEE Transactions on Automatic Control (\textbf{TAC}), pp. 1--8, 2026}
        {\plink{Paper}{https://arxiv.org/abs/2404.00481}
         \plink{Slides}{https://drive.google.com/file/d/1vuGl7mLdf4Muuaur2buQaIbgiV8pHmTu/view}}

\pubitem{\me, Chang Liu, Zhiqian Lan, Shengbo Eben Li, Wei Pan \& Angelo Alessandri}
        {Robust Bayesian Inference for Moving Horizon Estimation}
        {\textbf{Automatica}, 173, 112108, 2025}
        {\plink{Paper}{https://www.sciencedirect.com/science/article/abs/pii/S0005109824006010}
         \plink{Code}{https://github.com/wenhancao/robust-Bayesian-inference-for-MHE.git}}

\pubitem{Jingliang Duan, \me, Yang Zheng \& Lin Zhao}
        {On the Optimization Landscape of Dynamic Output Feedback Linear Quadratic Control}
        {IEEE Transactions on Automatic Control (\textbf{TAC}), 69(2):920--935, 2024}
        {\plink{Paper}{https://ieeexplore.ieee.org/abstract/document/10124022}
         \plink{Code}{https://github.com/soc-ucsd/LQG_gradient/tree/master/dLQR}}

\pubitem{Kangjie Zhou, Zhejia Wen, Zhiyong Zhuo, Zike Yan, Pengying Wu, Shuaiyang Li, Han Gao, Kang Ding, \me, Wei Pan \& Chang Liu}
        {CoINS: Counterfactual Interactive Navigation via Skill-Aware VLM}
        {IEEE Transactions on Automation Science and Engineering (\textbf{T-ASE}), 2026}
        {\plink{Paper}{https://arxiv.org/pdf/2601.03956}
         \plink{Project}{https://coins-internav.github.io/}}

\pubitem{Tianyi Zhang, \me, Chang Liu, Tao Zhang, Jiangtao Li \& Shengbo Eben Li}
        {Robust State Estimation for Legged Robots with Dual Beta Kalman Filter}
        {IEEE Robotics and Automation Letters (\textbf{RA-L}), 10(8):7955--7962, 2025}
        {\plink{Paper}{https://arxiv.org/abs/2411.11483}}

\pubitem{Shiqi Liu, \me, Chang Liu, Tianyi Zhang \& Shengbo Eben Li}
        {Convolutional Unscented Kalman Filter for Multi-Object Tracking with Outliers}
        {IEEE Transactions on Intelligent Vehicles (\textbf{T-IV}), 2024}
        {\plink{Paper}{https://ieeexplore.ieee.org/document/10643343}}

\pubitem{Tianyi Zhang, \me, Chang Liu, Feihong Zhang, Wei Wu \& Shengbo Eben Li}
        {NANO-SLAM: Natural Gradient Gaussian Approximation for Vehicle SLAM}
        {In IEEE International Conference on Intelligent Transportation Systems (\textbf{ITSC}), pp. 4014--4019, 2025}
        {\plink{Paper}{https://arxiv.org/abs/2504.19195}}

\pubitem{Tianyi Zhang, \me\ \& Shengbo Eben Li}
        {Natural Gradient Gaussian Approximation Filter with Positive Definiteness Guarantee}
        {In American Control Conference (\textbf{ACC}), pp. 3224--3229, 2026}
        {\relax}

\pubitem{\me, Chang Liu, Zhiqian Lan, Yingxi Piao \& Shengbo Eben Li}
        {Generalized Moving Horizon Estimation for Nonlinear Systems with Robustness to Measurement Outliers}
        {In American Control Conference (\textbf{ACC}), 2023}
        {\plink{Paper}{https://ieeexplore.ieee.org/document/10156391}
         \plink{Code}{https://github.com/wenhancao/robust-Bayesian-inference-for-MHE.git}
         \plink{Slides}{https://drive.google.com/file/d/1I4Ihmdwlf3kf3sYs4YXGwT6WWOHkr-Pd/view?usp=sharing}}

\pubitem{\me, Jingliang Duan, Shengbo Eben Li, Chen Chen, Chang Liu \& Yu Wang}
        {Primal-Dual Estimator Learning Method with Feasibility and Near-Optimality Guarantees}
        {In IEEE Conference on Decision and Control (\textbf{CDC}), 2022}
        {\plink{Paper}{https://ieeexplore.ieee.org/document/9992814}
         \plink{Slides}{https://drive.google.com/file/d/1Nueg0BkHXsr6s7BwFkDdz3MGJXMRvt9i/view?usp=sharing}}

\pubitem{\me, Jianyu Chen, Jingliang Duan, Shengbo Eben Li, Yao Lyu, Ziqing Gu \& Yuhang Zhang}
        {Reinforced Optimal Estimator}
        {In Modeling, Estimation and Control Conference (\textbf{MECC}), 2021. \textbf{(Student Best Paper Finalist)}}
        {\plink{Paper}{https://www.sciencedirect.com/science/article/pii/S240589632102245X}
         \plink{Slides}{https://drive.google.com/file/d/1ZXlsONIwnuPpb4IFkeCRYL7clDx7hTor/view?usp=sharing}}

\end{publist}


\section{预印本与在审论文}
\begin{publist}[{[P\arabic*]}]

\pubitem{\me, Keyu Yan, Yang Feng, Rick Siow Mong Goh \& Jun Zhou}
        {HIRE: Hippocampal Replay for Test-Time Agentic Reinforcement Learning}
        {\iclrstatus}
        {}

\pubitem{Keyu Yan$^*$, \mestar, Shen'ao Wang, Yujie Yang, Wu Junke, Xuyang Chen \& Lin Zhao}
        {Generating Parallel Worlds of a Flight}
        {\iclrstatus}
        {}

\pubitem{Keyu Yan, Xuyang Chen, \me\ \& Lin Zhao}
        {Evaluating OOD Actions in Offline RL via Quantile Advantage}
        {\iclrstatus}
        {\plink{Code}{https://anonymous.4open.science/r/adac-14D0}}

\pubitem{\me, Keyu Yan, Kuangyi Wang, Xuyang Chen \& Lin Zhao}
        {Trajectory Editing as Posterior Inference for Offline Reinforcement Learning}
        {\iclrstatus}
        {\plink{Code}{https://anonymous.4open.science/r/bite-055D}}

\pubitem{Xuyang Chen, Keyu Yan, Danze Chen, Hongzong Li, \me\ \& Lin Zhao}
        {Learning Robotic Skills from Offline Preference-based RL}
        {\iclrstatus}
        {\plink{Code}{https://anonymous.4open.science/r/PSAC-6686}}

\pubitem{Rui Huang, Xin Chen, \me, Lidong Li, Yizhe Xiao, Yichao Gao, Yuqi Shi \& Lin Zhao}
        {Agile Modular Aerial Robot Load Transport with Arbitrary Configurations}
        {\tmecstatus}
        {\relax}

\pubitem{Han Gao, \me, Ieng Hou U, Kangjie Zhou, Ying Sun, Angelo Alessandri \& Chang Liu}
        {Chance-Constrained Gaussian Filtering}
        {\automaticastatus}
        {\relax}

\end{publist}


\section{荣誉与奖励}
\award{Eric and Wendy Schmidt AI in Science 博士后奖学金（每年遴选不超过五人）}{2026}
\award{清华大学海外访学基金}{2022}
\award{MECC 学生最佳论文奖入围}{2021}
\award{国家奖学金}{2016}
\award{北京交通大学一等奖学金}{2016、2017、2018}

\section{开源软件}
\textbf{General Optimal control Problems Solver (GOPS)}：面向工业应用、用于构建实时高性能控制器的强化学习求解器软件包。我主要负责 \textbf{trainer}、\textbf{sampler} 和 \textbf{buffer} 模块的核心设计与实现。
\plink{文档}{https://gops.readthedocs.io/en/latest/index.html}
\plink{代码}{https://github.com/Intelligent-Driving-Laboratory/GOPS}
\plink{论文}{https://www.sciencedirect.com/science/article/pii/S2772424723000070}

\section{代表性科研项目}
\textbf{MindGPT 的强化学习后训练} \hfill 2024 年 11 月\par
项目成员，\href{https://www.lixiang.com/}{理想汽车}实习
\vspace{\gapM}

\textbf{开放道路自动驾驶决策与控制的强化学习方法} \hfill 2024 年 5 月\par
项目成员；\href{https://global.toyota/en/}{丰田}与\href{https://www.idriverplus.com/}{智行者}公司支持
\vspace{\gapM}

\textbf{仓储自动物流车辆的状态估计} \hfill 2022 年 5 月\par
学生项目负责人；\href{https://www.geekplus.com/en/}{极智嘉（Geek+）}支持，相关技术已应用于公司产品
\vspace{\gapM}

\textbf{网联自动驾驶电动车辆的网络化建模与协同控制} \hfill 2020 年 11 月\par
学生项目负责人；科技部支持

\section{受邀报告}
\begin{itemize}[leftmargin=0pt,itemindent=0pt,labelsep=0pt,label={},itemsep=\gapS,parsep=0pt,topsep=3pt plus 1pt]
\justifying
\setlength{\parindent}{0pt}
\interlinepenalty=10000\relax
\item \textit{NANO Filter: Bayesian Filtering with Natural Gradient Gaussian Approximation}\par
清华大学天文系，中国北京；蔡峥教授邀请，2024 年 8 月
\item \textit{Convolutional Bayesian Filtering}\par
清华大学数学科学系，中国北京；丘成栋教授邀请，2024 年 2 月
\item \textit{Learning-based State Estimation Methods}\par
慕尼黑工业大学，德国慕尼黑（线上）；Sandra Hirche 教授邀请，2023 年 2 月
\end{itemize}

\section{教学经历}
\textbf{清华大学车辆与运载学院} \hfill 2022 年 2 月 -- 2022 年 5 月\par
课程助教：150051 智能汽车（Intelligent Vehicles）

\section{学术服务}
\textbf{会议审稿：} NeurIPS、ICML、ICLR、AAMAS、L4DC、CDC、ACC、IFAC NMPC\par
\vspace{2pt}
\textbf{期刊审稿：} TPAMI、T-RO、T-RL、TAC、Automatica、RA-L、T-ITS、T-ASE、TII、TNNLS
\end{document}
